% Figure: the concentric-cylinder (or sphere-in-enclosure) radiation
% geometry used to deepen the two-surface exchange algebra. An inner
% convex surface (1) sees only the outer surface (2), so F12 = 1; the
% outer surface sees both the inner body and itself, so F22 = 1 - A1/A2.
% Drawn as two nested circles with rays to make the self-view explicit.
\begin{figure}[htbp]
\centering
\begin{tikzpicture}[scale=1.0]
% outer enclosure
\draw[engblue, very thick] (0,0) circle (2.6);
% inner convex body
\draw[engred, very thick, fill=engred!10] (0,0) circle (1.0);
\node[font=\small, engred] at (0,0) {$A_1,\ T_1$};
\node[font=\small, engblue] at (0,2.05) {$A_2,\ T_2$};
% a ray leaving the inner body, reaching the outer surface (F12 = 1)
\draw[engblack, ->, thick] (0.71,0.71) -- (1.84,1.84);
\node[font=\scriptsize, engblack, anchor=west] at (1.5,1.3) {$F_{12}=1$};
% a ray leaving the outer surface and returning to the outer surface
% (the self-view F22), shown as an arc-chord that misses the inner body
\draw[enggray, ->, thick] (-2.45,0.86) to[bend right=35] (-0.86,-2.45);
\node[font=\scriptsize, enggray, anchor=east] at (-2.55,-1.0)
  {$F_{22}=1-\dfrac{A_1}{A_2}$};
% separation gap label
\draw[englime!50!black, <->] (1.0,-0.0) -- (2.6,-0.0);
\node[font=\scriptsize, englime!50!black, anchor=north] at (1.8,-0.1)
  {evacuated gap};
\end{tikzpicture}
\caption[Concentric-body radiation geometry and its view
factors]{The radiation geometry of a convex body inside an enclosure,
the case the two-surface exchange algebra reduces to most often. Every
ray leaving the inner convex surface lands on the enclosure, so the view
factor from the inner body to the enclosure is exactly one. The enclosure
sees both the inner body and itself, and the summation rule fixes its
self-view factor at $F_{22}=1-A_1/A_2$ once reciprocity gives
$F_{21}=A_1/A_2$. Substituting these into the general two-surface formula
collapses it to a single expression in the area ratio and the two
emissivities, the working form for a tank in a chamber, a satellite in
space, and a wire in a duct.}
\label{fig:vol04ch05:concentric-cylinders}
\end{figure}
